% Sections 12.1 to 12.6, translated from sv/chapters/12/theory.tex.

\section{Rate of change and the difference quotient}
\label{c12:sec:andringskvot}
The \textbf{difference quotient} (the average rate of change) of $f(x)$ over the interval
$[x, x + h]$ is
\[
  \frac{f(x + h) - f(x)}{h}.
\]
Geometrically, this is the slope of the \textbf{secant} through the points $(x, f(x))$ and
$(x + h, f(x + h))$.

\section{The definition of the derivative}
\label{c12:sec:definition}
The derivative $f'(x)$ is the limit of the difference quotient as $h \to 0$:
\[
  f'(x) = \lim_{h \to 0} \frac{f(x + h) - f(x)}{h}
\]
Geometrically, $f'(x)$ is the \textbf{slope of the tangent} to the curve $y = f(x)$ at the point
$x$. The derivative is also called the \textbf{instantaneous rate of change}.

\section{Rules of differentiation for polynomials}
\label{c12:sec:regler}

\begin{narrowtable}{@{}ll@{}}
  \thead{Function} & \thead{Derivative} \\
  \midrule
  $f(x) = C$ (constant) & $f'(x) = 0$ \\
  $f(x) = x^n$ & $f'(x) = nx^{n-1}$ \\
  $f(x) = Cx^n$ & $f'(x) = Cnx^{n-1}$ \\
  $f(x) = u(x) + v(x)$ & $f'(x) = u'(x) + v'(x)$ \\
\end{narrowtable}

\begin{exempel}
\[
  f(x) = 3x^4 - 2x^2 + 5 \quad \Rightarrow \quad f'(x) = 12x^3 - 4x
\]
\end{exempel}

\section{The equation of the tangent}
\label{c12:sec:tangent}
The tangent to $y = f(x)$ at the point $(x_0, f(x_0))$ has the equation
\[
  y = f'(x_0)(x - x_0) + f(x_0).
\]

\section{Turning points -- maximum and minimum}
\label{c12:sec:extrempunkter}
A \textbf{stationary point} satisfies $f'(x_0) = 0$. Whether it is a maximum or a minimum is
decided with the second-derivative test:
\begin{itemize}
  \item $f''(x_0) > 0$: \textbf{minimum} (the curve is concave up),
  \item $f''(x_0) < 0$: \textbf{maximum} (the curve is concave down).
\end{itemize}

\section{Worked examples}
\label{c12:sec:typexempel}

\begin{exempel}[Differentiating polynomial functions]
Differentiate
\begin{align*}
  \textbf{a)}\quad f(x) &= -2x^2 + 2x + 4, \\
  \textbf{b)}\quad f(x) &= 3x^3 - 6x^2 + \frac{3x}{4} - 5, \\
  \textbf{c)}\quad f(x) &= -x^4 + x^3 + \frac{2x^2}{3} - 3x + 2.
\end{align*}
The rule for $Cx^n$ is applied term by term:
\begin{align*}
  \textbf{a)}\quad f'(x) &= -4x + 2 \\
  \textbf{b)}\quad f'(x) &= 9x^2 - 12x + \frac{3}{4} \\
  \textbf{c)}\quad f'(x) &= -4x^3 + 3x^2 + \frac{4x}{3} - 3
\end{align*}
\end{exempel}

\begin{exempel}[Analysing a parabola]
The function is $f(x) = -x^2 + 6x - 5$.
\par\noindent\textbf{a)} Find $f'(x)$.
\[
  f'(x) = -2x + 6
\]
\textbf{b)} Solve $f'(x) = 0$.
\[
  -2x + 6 = 0 \quad \Rightarrow \quad x = 3
\]
\textbf{c)} Determine the type with $f''(x)$.
\[
  f''(x) = -2 < 0 \quad \Rightarrow \quad \text{maximum}
\]
\textbf{d)} Calculate the maximum value.
\[
  f(3) = -9 + 18 - 5 = 4
\]
The maximum point is $(3, 4)$.
\end{exempel}

\begin{exempel}[Current in an RC circuit]
The current is approximated by $i(t) = -0.4t^3 + 2.4t^2 - 3t + 1.2$ amperes, with the time $t$ in
seconds.
\par\noindent\textbf{a)} Calculate $i'(t)$.
\[
  i'(t) = -1.2t^2 + 4.8t - 3
\]
\textbf{b)} Solve $i'(t) = 0$ with the PQ-formula.
\begin{gather*}
  t^2 - 4t + 2.5 = 0 \quad \Rightarrow \quad t = 2 \pm \sqrt{1.5} \\
  t_1 \approx 0.78\,\text{s}, \qquad t_2 \approx 3.22\,\text{s}
\end{gather*}
\textbf{c)} The second derivative is $i''(t) = -2.4t + 4.8$, and
\begin{alignat*}{2}
  i''(0.78) &\approx 2.9 > 0 &\qquad& \Rightarrow \quad \text{minimum at } t_1, \\
  i''(3.22) &\approx -2.9 < 0 && \Rightarrow \quad \text{maximum at } t_2.
\end{alignat*}
\textbf{d)} The current at the stationary points is
\[
  i(0.78) \approx 0.13\,\text{A} \ \text{(minimum)}, \qquad
  i(3.22) \approx 3.07\,\text{A} \ \text{(maximum)}.
\]
\end{exempel}
